\subsection{Esempio: lo spin di una particella in 3 dimensioni}
L'esempio più importante di spazio di Hilbert prodotto è dato dal grado di libertà di spin,
che, nella meccanica quantistica non relativistica, deve essere postulato (mentre emerge dalla
formulazione di Dirac).

$\vec{S}$ è un operatore vettoriale che segue l'algebra del momento angolare
$[S_i,S_j] = i\hbar\epsilon_{ijk}S_k$; dalla teoria del momento angolare si trova che, scelto il SCOC
$\{S^2,S_z\}$, gli autovalori degli autostati simultanei sono
$$
\begin{dcases}
    S^2\ket{s,s_z} = \hbar^2s(s+1)\ket{s,s_z} & s = \frac{n}{2}, n\in\mathbb{Z} \\
    S_z\ket{s,s_z} = \hbar s_z\ket{s,s_z} & s_z = -s,-s+1,...,0,...,s-1,s\\
\end{dcases}
$$
lo spazio di Hilbert dei gradi di libertà di spin è dunque uno spazio vettoriale di dimensione
$2s+1$ basato sul campo complesso: $\mathcal{H}_\text{spin} = \mathbb{C}^{2s+1}$.

Visto che questi operatori non "vedono" la parte spaziale, lo spazio di Hilbert totale di una
particella con spin $s$ in 3 dimensioni è
$$ \mathcal{H} = L^2(\mathbb{R}^3) \otimes \mathbb{C}^{2s+1}$$
Scelta come base $\{ \ket{\vec{x},s_z} = \ket{\vec{x}}\otimes\ket{s_z} \}$\footnote{con
$\ket{s_z = s} = e_1$ e $\ket{s_z=-s} = e_{2s+1}$ elementi della base standard di $\mathbb{C}^{2s+1}$},
si può scrivere il generico stato del sistema
$$\ket{\psi} = \sum_{s_z=-s}^{s}\int_{}^{}d^3\vec{x} \ket{\vec{x},s_z}\bra{\vec{x},s_z}\ket{\psi}
:= \sum_{s_z=-s}^{s}\int_{}^{}d^3\vec{x} \psi_{s_z}(\vec{x})\ket{\vec{x},s_z}$$
dunque
$$ \bra{\vec{x}}\ket{\psi}=\psi(\vec{x}) = \sum_{s_z=-s}^{s}\psi_{s_z}(\vec{x})\ket{s_z} =
    \begin{pmatrix}
	\psi_s(\vec{x}) \\
	\vdots \\
	\psi_{-s}(\vec{x})
    \end{pmatrix}$$
che è detta funzione d'onda vettoriale.

\section{Postulato 1 della Meccanica Quantistica}
\noindent\fbox{%
    \parbox{\textwidth}{%
    \textbf{I possibili risultati della misura di un osservabile F sono solo i suoi autovalori}
    $\{\mathbf{f_i}\}$\textbf{. Ciascun autovalore ha una probabilità}
    $\mathbf{\mathcal{P}_{\psi}(f_i)}$ \textbf{dipendente dallo stato del sistema.}
    }%
}

Questo è detto, per ovvii motivi, postulato della misura.

Per definire queste probablità bisogna distinguere 3 casi:

\begin{itemize}
    \item $\mathbf{f_i \in Spd(F)}$ \textbf{non degenere}
    
	Dato l'autostato (normalizzato) $\ket{f_i}$ e uno stato $\ket{\psi}$, si \textbf{definisce}	
	$$\mathcal{P}_{\psi}(f_i):= \left| \bra{f_i}\ket{\psi} \right|^2 =
	\bra{\psi}\ket{f_i}\bra{f_i}\ket{\psi} := \bra{\psi}P_{f_i}\ket{\psi}$$
	dove $P_{f_i} := \ket{f_i}\bra{f_i}$ è il \textbf{proiettore} sullo stato $\ket{f_i}$

    \item $\mathbf{f_i \in Spd(F)}$ \textbf{degenere}
	Si supponga che l'autovalore $f_i$ abbia una degenerazione di grado $n_i$:
	si può sempre costruire una base per l'autospazio degenere $\left\{\ket{f_i,\alpha_\mu}\right\}$,
	con $\mu = 1,2...,n_i$;

	Si \textbf{definisce} dunque
	$$\mathcal{P}_\psi(f_i):=\sum_{\mu=1}^{n_i}\left|\bra{f_i,\alpha_\mu}\ket{\psi}\right|^2
	= \bra{\psi}\left(\sum_{\mu=1}^{n_i}\ket{f_i,\alpha_\mu}\bra{f_i,\alpha_\mu} \right)\ket{\psi}:=
	\bra{\psi}P_{f_i}\ket{\psi}$$
	avendo opportunamente ridefinito il proiettore.
    \item $\mathbf{f\in Spc(F)}$

	Nel caso di spettro continuo, si definisce una \textbf{densità di probabilità} nel modo seguente:
	$$\mathcal{P}_\psi(f,f+\Delta f) := \int_{f}^{f+\Delta f}df \bra{\psi}P_f\ket{\psi}$$
	dove $P_f:=\ket{f}\bra{f}$ (nel caso non degenere).	
\end{itemize}

\subsection{Valor medio di un'osservabile}
Data un'osservabile $F$ e uno stato normalizzato $\ket{\psi}$, il suo valor medio è il numero
$$\langle F \rangle_\psi = \bra{\psi}F\ket{\psi} := \sum_{i}^{}f_i\mathcal{P}_\psi(f_i) =
\sum_{i}^{}f_i\left| \bra{f_i}\ket{\psi} \right|^2$$
un modo più generico di scrivere la stessa quantità, come sarà chiaro tra poco, è il seguente:

$$
\bra{\psi}F\ket{\psi} = \sum_{i}^{}\bra{\psi}\ket{i}\bra{i}F\ket{\psi} =
\sum_{i}^{}\bra{i}F\ket{\psi}\bra{\psi}\ket{i} = Tr(FP_{\psi})
$$
con $\left\{\ket{i}\right\}$ base ONC qualunque, in quanto la traccia di un operatore è invariante per
cambio di base.

Similmente si può ricavare che $\mathcal{P}_{\psi}(f_i) = Tr(P_{f_i}P_{\psi})$.

\subsection{Stati misti e matrice densità}
Prima di procedere con la trattazione del postulato della misura, è necessario introdurre il concetto
di \textbf{stato misto}.

\textbf{Def.} Uno stato misto è una miscela di stati quantistici puri $\ket{\psi_k}$ tale che la
probabilità classica di trovare uno degli stati $\ket{\psi_k}$ sia $\omega_k$, con
$$
\sum_{k=1}^{N}\omega_k=1.
$$
Si consideri ora un'osservabile $F$: il valore medio di $F$ sarà la media classica dei valori medii:
$$ \sum_{k=1}^{N} \omega_k \bra{\psi_k}F\ket{\psi_k} = \sum_{k=1}^{N}\omega_k Tr(FP_k) =
Tr\left(F \sum_{k=1}^{N} \omega_k \ket{\psi_k}\bra{\psi_k}\right) := Tr(F\rho):=\langle F\rangle_\rho$$

\textbf{Def.} L'operatore $\rho :=\sum_{k=1}^{N}\omega_k\ket{\psi_k}\bra{\psi_k}$ si dice \textbf{matrice
densità} e rappresenta lo stato misto.

\textit{Oss.} La matrice densità di uno stato puro è il caso banale $\rho_\psi = \ket{\psi}\bra{\psi}$:
questo rende più chiara la notazione $\langle F\rangle_\psi = Tr(FP_\psi) = Tr(F\rho_\psi)$

\subsubsection{Proprietà della matrice densità}
\begin{enumerate}
    \item Hermiticità $\rho = \rho^{\dagger}$;
    \item $Tr(\rho) = \sum_{k}^{}\omega_k Tr(P_k) =
	\sum_{k}^{}\omega_k\sum_{i}^{}\bra{i}\ket{\psi_k}\bra{\psi_k}\ket{i} = \sum_{k}^{}\omega_k  = 1$;
    \item Autovettori e autovalori $\rho\ket{\psi_l}
	= \sum_{k}^{}\omega_k\bra{\psi_k}\ket{\psi_l}\ket{\psi_k} = \omega_l\ket{\psi_l}$, che è
	ovvio, in quanto $\rho$ è diagonale nella base $\left\{\ket{\psi_k}\right\}$;
    \item $\forall \ket{v}\in\mathcal{H}, \bra{v}\rho\ket{v}
	= \sum_{k}^{} \omega_k \bra{v}\ket{\psi_k}\bra{\psi_k}\ket{v}
	= \sum_{k}^{} \omega_k \left|\bra{v}\ket{\psi_k} \right|^2 \geq 0$;
    \item Idempotenza $\rho^2 = \rho$, quindi $Tr(\rho^2) = Tr(\rho) = 1$.
\end{enumerate}

I risultati della misura dell'osservabile $F$ sono sempre gli autovalori $f_i$, pesati con la probablità
sia classica $\omega_k$ sia quantistica $\mathcal{P}_{\psi_k}(f_i)$:
$$\mathcal{P}_\rho(f_i) := \langle P_{f_i}\rangle_\rho = Tr(P_{f_i}\rho) = 
\sum_{k=1}^{N}\omega_k \mathcal{P}_{\psi_k}(f_i)\in\mathbb{R}$$
\textbf{Esempio} Quale è la differenza tra uno stato puro e uno stato misto?

Si consideri un operatore $A$ con spettro discreto $\left\{a_i\right\}$ e due stati quantistici
$$\ket{\psi} = \sum_{i}^{}a_i\ket{a_i} \text{ e } \rho = \sum_{i}^{}\omega_i\ket{a_i}\bra{a_i}
\text{, con } \omega_i = |c_i|^2
$$
Dalla misura di $A$ si trovano risultati identici:
$$\mathcal{P}_\psi(a_i) = \bra{\psi}P_{a_i}\ket{\psi} = \left|\bra{a_i}\ket{\psi}\right|^2 =
|c_i|^2 = \omega_i$$
$$\mathcal{P}_\rho(a_i) = Tr(P_{a_i}\rho) =
Tr\left( \sum_{j}^{}\omega_j \left( \ket{a_i}\bra{a_i}\right)\left(\ket{a_j}\bra{a_j}\right) \right) = 
\sum_{j}^{}\omega_j \bra{a_j}\ket{a_i}\bra{a_i}\ket{a_j} = \omega_i
$$
Se però si considera un operatore $B$ tale che $\left[A,B\right] \neq0$, le probabilità di misurare un suo
autovalore $b$ sono:
$$\mathcal{P}_\rho(b) = Tr(P_b\rho) = Tr\left( \sum_{k}^{}\omega_k\ket{a_k}\bra{a_k}\ket{b}\bra{b}\right)=
\sum_{k}^{}\omega_k\left|\bra{a_k}\ket{b}\right|^2
$$
mentre
\begin{align*}
    \mathcal{P}_\psi(b)&=\bra{\psi}P_b\ket{\psi}=
    \sum_{i}^{}\sum_{j}c_i^*c_j\bra{a_i}\ket{b}\bra{b}\ket{a_j} =
    \sum_{i=j}\left[...\right] +\sum_{i>j}\left[...\right]+\sum_{i<j}\left[...\right] = \\
    &=\sum_{i}|c_i|^2\left|\bra{a_i}\ket{b}\right|^2 +
    \sum_{i>j}c_i^*c_j\bra{a_i}\ket{b}\bra{b}\ket{a_j} + c_j^*c_i\bra{a_j}\ket{b}\bra{b}\ket{a_i}=\\
    &=\sum_{i}\omega_i\left|\bra{a_i}\ket{b}\right|^2 +
    2\Re\left( \sum_{i>j} c_i^*c_j\bra{a_i}\ket{b}\bra{b}\ket{a_j}\right) 
\end{align*}

La differenza tra queste probabilità può essere intesa nel seguente modo: la miscela di stati puri $\rho$
non comprende le fasi dei coefficienti di probabilità $\arg(c_i)$, mentre lo stato puro $\psi$ sì; nella
probabilità associata a $\rho$ c'è solo il prodotto tra la probabilità classica e quella quantistica, 
in quella associata a $\psi$ c'è anche un termine dato dall'interferenza tra le fasi dei coefficienti
$c_i$.

Un altro modo per vedere questa differenza è considerare la matrice densità associata a $\psi$
$\rho_\psi = \ket{\psi}\bra{\psi}$: essa non è diagonale nella base di autostati di $A$, infatti
$$\rho_\psi = \ket{\psi}\bra{\psi} = \sum_{i}\sum_{j}^{}c_ic_j^*\ket{a_i}\bra{a_j}$$
mentre la matrice $\rho = \sum_{k}\left|c_k\right|^2\ket{a_k}\bra{a_k}$ lo è, quindi in $\rho_\psi$ sono
contenute più informazioni, cioè le fasi dei $c_i$.

\textbf{Oss.} Gli elementi diagonali di una matrice densità $\rho$ si chiamano \textit{popolazioni}, 
mentre quelli non diagonali si chiamano \textit{coerenze}, dunque $\rho$ si dice
\textbf{sovrapposizione coerente} di stati, mentre $\rho_\psi$ si dice
\textbf{sovrapposizione incoerente} di stati.

